\nuevaReceta{Fajitas de Pollo y Tres Pimientos}{4 personas}{35 min}{receta:fajitas}
    
    \begin{minipage}[t]{0.35\textwidth}
        \section*{Ingredientes}
        \begin{itemize}[leftmargin=*]
            \raggedright
            \item 600 g de pechuga de pollo
            \item 3 pimientos (verde, rojo y amarillo)
            \item 2 chalotas o 1 cebolla
            \item 1 sobre de sazonador de fajitas
            \item sazonador barbacoa y de carne
            \item hierbas provenzales
            \item 8 tortillas de trigo
            \item \textbf{Guacamole:} (Ver receta en pág. \pageref{receta:guacamole})
        \end{itemize}
    \end{minipage}
    \hfill
    \begin{minipage}[t]{0.6\textwidth}
        \section*{Preparación}
        \begin{enumerate}
            \item \textbf{Sazonar el pollo:} Cortar la carne en trozos medianos/pequeños. Sazonar con barbacoa, carne, hierbas provenzales y un toque del sobre de fajitas.
            \item \textbf{Verduras:} Cortar los pimientos en tiras y luego por la mitad. Picar la chalota de forma similar.
            \item \textbf{Sofrito:} En una sartén con aceite, dorar la cebolla. Añadir los pimientos con un poco de sazonador barbacoa y de fajitas. Cocinar 10-15 min.
            \item \textbf{Cocción lenta:} Añadir el pollo y más sazonador de fajitas. Tapar la sartén y bajar el fuego para que se cocine en su propio jugo, removiendo ocasionalmente.
            \item \textbf{Toque final:} Destapar y subir el fuego al final para consumir el caldo al gusto. Calentar las tortillas y servir con una base de guacamole.
        \end{enumerate}
    \end{minipage}

    \vspace{1em}
    \begin{tcolorbox}[colback=orange!5, colframe=orange!50, title=El rincón del Chef]
        \begin{minipage}{0.25\textwidth}
            \centering
            \usebox{\miqrMexicano} % <--- Aquí llamamos al QR ya renderizado
        \end{minipage}
        \begin{minipage}{0.7\textwidth}
            \textbf{\faMusic\ Música:} Escanea para cocinar con ritmo. \\
            \textbf{\faLightbulb\ Nota:} \textit{Unta el guacamole en la tortilla caliente.}
        \end{minipage}
    \end{tcolorbox}
\end{receta}
